\section{结论}

加密矿企转型 AI 数据中心，本质上是一场“算力迁移”。在这场迁移中，税收政策是影响项目选址、投资时点和商业模式设计的重要变量之一。

矿企需要考虑的不仅是哪个州提供的优惠最多，还要评估项目能否取得优惠资格、优惠可以持续多久，以及取得资格后能否持续满足投资、就业、报告和记录保存要求。申请过晚，或在取得所需免税认证之前进行采购，可能使本可获得的节税收益落空。后续运营中未能履行绩效承诺，可能导致补缴税款、支付利息，甚至失去优惠资格。

美国各州对数据中心的政策态度正处于微妙的转折点。一些辖区继续提供力度较大的税收优惠，另一些则暂停新项目审批或缩小既有优惠范围。随着政策窗口收窄，“先行动、后申请”的做法正变得越来越难以奏效。对于计划转型的加密矿企，税收优惠申请及长期税务规划需要与投资决策同步推进。

\FloatBarrier
